\section{Processo di Sviluppo Incrementale organizzato in Sprint}
\label{sec:processo_sviluppo_sprint}

\subsection{Panoramica del Ciclo di Sviluppo Iterativo}
In aderenza alla metodologia dello \emph{Specification-Driven Development} (SDD), lo sviluppo di TEPO si è articolato in quattro sprint di durata bisettimanale. Ciascun ciclo ha trasformato una porzione del Product Backlog in codice sorgente collaudato, documentando contestualmente l'evoluzione dell'architettura logica e della struttura delle classi.
La Figura~\ref{fig:roadmap_sprint} riassume la traiettoria temporale e i deliverable di incremento prodotti dal team.

\begin{figure}[H]
    \centering
    \begin{tikzpicture}[
        sprintbox/.style={
            draw=tepoblue, fill=tepoblue!6, rounded corners=4pt, thick,
            text width=6.7cm, minimum height=3.6cm, align=left, font=\scriptsize, inner sep=6pt
        },
        titlebox/.style={
            fill=tepoblue!25, draw=tepoblue, rounded corners=2pt,
            font=\small\bfseries, align=center, text width=6.4cm, minimum height=0.55cm
        },
        arrow/.style={-Latex, very thick, draw=tepoblue!80}
    ]
        % Row 1: Sprint 1 and Sprint 2
        \node[sprintbox] (s1) at (0,0) {
            \\[0.35cm]
            \textbf{Focus:} Dominio, Persistenza e Shell\\
            $\bullet$ US-01: Catalogo Edifici Comunali\\
            $\bullet$ US-02: Catalogo Interventi e Tipologie\\
            $\bullet$ US-08: Cartografia Interattiva Leaflet\\
            $\bullet$ US-09: Dashboard di Controllo e Navigazione\\[3pt]
            \textbf{Incremento:} Schema Flyway V1, Entità JPA, Shell React\\
            \textbf{Complessità:} 16 Story Points
        };
        \node[titlebox] at (s1.north) {Sprint 1: Fondamenta e Catalogo};

        \node[sprintbox] (s2) at (7.8cm, 0) {
            \\[0.35cm]
            \textbf{Focus:} Ottimizzazione Esatta e Spiegazioni\\
            $\bullet$ US-03: Scenario Builder con Vincoli\\
            $\bullet$ US-04: Risoluzione Ottima OR-Tools CP-SAT\\
            $\bullet$ US-05: Motore Spiegazioni Deterministiche\\
            $\bullet$ US-14: Audit Trail e Hash Canonico SHA-256\\[3pt]
            \textbf{Incremento:} Schema V2, Core CP-SAT, Tie-Break a 4 Passi\\
            \textbf{Complessità:} 34 Story Points
        };
        \node[titlebox] at (s2.north) {Sprint 2: Algoritmo e Solutore};

        % Row 2: Sprint 3 and Sprint 4
        \node[sprintbox] (s3) at (0, -4.9cm) {
            \\[0.35cm]
            \textbf{Focus:} Analisi Comparativa e Dati Esterni\\
            $\bullet$ US-06: Confronto Scenari What-If (Slot A/B)\\
            $\bullet$ US-07: Analisi di Sensibilità sul Budget (Sweep)\\
            $\bullet$ US-10: Importazione Guidata CSV con Dry-Run\\[3pt]
            \textbf{Incremento:} Sensitivity Engine, Compare View, Store Client\\
            \textbf{Complessità:} 21 Story Points
        };
        \node[titlebox] at (s3.north) {Sprint 3: What-If e Sensibilità};

        \node[sprintbox] (s4) at (7.8cm, -4.9cm) {
            \\[0.35cm]
            \textbf{Focus:} Assistente AI e Rilascio in Esercizio\\
            $\bullet$ US-11: Assistente AI Human-in-the-Loop\\
            $\bullet$ US-12: Fallback Rule-Based Locale\\
            $\bullet$ US-13: Accessibilità WCAG AA (Relazione / Sala)\\
            $\bullet$ US-15: Esportazione Relazione e Stampa\\[3pt]
            \textbf{Incremento:} AI Ibrida, Container Docker, Deploy Render\\
            \textbf{Complessità:} 16 Story Points
        };
        \node[titlebox] at (s4.north) {Sprint 4: AI e Deployment Cloud};

        % Flow arrows connecting the 4 sprints
        \draw[arrow] (s1.east) -- (s2.west);
        \draw[arrow] ([xshift=-1.5cm]s2.south) -- ++(0, -0.45cm)
            -- node[pos=0.30, fill=white, draw=tepoblue!70, rounded corners=2pt, inner sep=2.5pt, font=\scriptsize\bfseries, text=tepoblue!90] {Incremento}
               ([xshift=1.5cm]s3.north |- 0, -2.6cm)
            -- ([xshift=1.5cm]s3.north);
        \draw[arrow] (s3.east) -- (s4.west);
    \end{tikzpicture}
    \caption{Roadmap evolutiva dei quattro sprint di sviluppo di TEPO.}
    \label{fig:roadmap_sprint}
\end{figure}

% -------------------------------------------------------------
% SPRINT 1
% -------------------------------------------------------------
\subsection{Sprint 1: Fondamenta di Dominio, Catalogo, Cartografia e Shell}

\subsubsection{Obiettivo dello Sprint e User Story Affrontate}
L'obiettivo primario dello Sprint 1 è stato l'allestimento dell'infrastruttura fondamentale del progetto, stabilendo il modello dei dati relazionale, popolando il dataset dimostrativo sintetico e sviluppando la shell applicativa con le funzionalità di navigazione, esplorazione tabellare e cartografica.
Le User Story completate sono:
\begin{itemize}[leftmargin=1.5em, itemsep=0.2em]
    \item \textbf{US-01 (Catalogo degli Edifici):} Visualizzazione tabellare interattiva dei 100 edifici comunali, con metriche di superficie, consumi annui kWh, spesa energetica in \euro{} e indicazione esplicita della provenienza sintetica (`SYNTHETIC`).
    \item \textbf{US-02 (Catalogo Interventi e Tipologie):} Esplorazione delle 300 opzioni di intervento suddivise per le 6 tipologie (Caldaie, Pompe di Calore, Isolamento, Serramenti, Fotovoltaico, Illuminazione LED), con relativi costi e benefici nominali.
    \item \textbf{US-08 (Cartografia Georeferenziata Leaflet):} Mappa interattiva del territorio di Trento, posizionamento degli edifici tramite coordinate geografiche (latitudine/longitudine) e popup informativo con riepilogo prestazionale.
    \item \textbf{US-09 (Dashboard di Controllo e Navigazione):} Shell unificata con navigazione laterale (`SideNav`), riepilogo sintetico del patrimonio edilizio e badge permanente per la dichiarazione di conformità demo.
\end{itemize}

\subsubsection{Piano di Implementazione Tecnico}
\begin{itemize}[leftmargin=1.4em, itemsep=0.2em]
    \item \textbf{Database:} Creazione della migrazione Flyway iniziale \texttt{V1\_\_init.sql}, con definizione delle tabelle \texttt{buildings}, \texttt{intervention\_types} e \texttt{intervention\_options}, vincoli di chiave primaria, chiave esterna e indice \texttt{idx\_options\_building}.
    \item \textbf{Backend:} Implementazione del package \texttt{it.unitn.tepo.domain} con le classi entità \texttt{Building}, \texttt{InterventionOption}, \texttt{InterventionType} e l'enumerazione \texttt{DataProvenance}. Realizzazione dei repository Spring Data JPA (\texttt{BuildingRepository}, \texttt{InterventionOptionRepository}, \texttt{InterventionTypeRepository}) e dei controller REST di consultazione (\texttt{BuildingController}, \texttt{InterventionController}, \texttt{MapController}).
    \item \textbf{Frontend:} Allestimento dell'ambiente Vite con TypeScript; sviluppo dei componenti di layout \texttt{AppShell}, \texttt{SideNav}, \texttt{StatusBand}; integrazione della libreria cartografica Leaflet con styling ad alto contrasto; realizzazione delle pagine \texttt{Dashboard.tsx}, \texttt{BuildingsPage.tsx}, \texttt{InterventionsPage.tsx} e \texttt{MapPage.tsx}.
\end{itemize}

\subsubsection{Evoluzione Architetturale dello Sprint 1}
La Figura~\ref{fig:cmp_sprint1} illustra il diagramma dei componenti dello Sprint 1, evidenziando la separazione tra il livello di presentazione React, i controller REST di lettura e il database relazionale.

\begin{figure}[H]
    \centering
    \begin{tikzpicture}[
        node distance=0.9cm and 1.2cm,
        comp/.style={draw=black!80, fill=tepoblue!10, rounded corners=3pt, thick, align=center, font=\scriptsize\bfseries, minimum height=0.9cm, minimum width=3.4cm},
        db/.style={cylinder, draw=black!80, fill=yellow!15, thick, shape border rotate=90, aspect=0.25, minimum width=2.4cm, minimum height=1.1cm, align=center, font=\scriptsize\bfseries},
        box/.style={draw=tepoblue!60, fill=tepoblue!3, rounded corners=4pt, dashed, inner sep=7pt},
        arrow/.style={-Latex, thick, draw=black!80}
    ]
        % Client Box
        \node[comp] (ui_shell) {AppShell / SideNav};
        \node[comp, right=0.6cm of ui_shell] (ui_views) {Dashboard / MapPage /\\BuildingsPage / InterventionsPage};
        \node[comp, below=0.6cm of ui_shell] (ui_api) {ApiClient (TypeScript)};
        \begin{scope}[on background layer]
            \node[box, fit=(ui_shell) (ui_views) (ui_api), label=above:{\scriptsize\textbf{Frontend SPA (React 18 / Vite)}}] (box_fe) {};
        \end{scope}

        % Backend Box
        \node[comp, below=1.5cm of ui_api] (ctrl) {BuildingController /\\InterventionController / MapController};
        \node[comp, right=0.8cm of ctrl] (repos) {BuildingRepository /\\InterventionOptionRepository};
        \begin{scope}[on background layer]
            \node[box, fit=(ctrl) (repos), label=above:{\scriptsize\textbf{Backend Services (Spring Boot 3.3)}}] (box_be) {};
        \end{scope}

        % Database
        \node[db, below=1cm of repos] (db_node) {PostgreSQL 16\\(Flyway V1)};

        % Connectors
        \draw[arrow] (ui_views) -- (ui_api);
        \draw[arrow] (ui_api) -- node[left, font=\tiny] {REST / JSON} (ctrl);
        \draw[arrow] (ctrl) -- (repos);
        \draw[arrow] (repos) -- node[right, font=\tiny] {JDBC / SQL} (db_node);
    \end{tikzpicture}
    \caption{Component Diagram allo stato dello Sprint 1.}
    \label{fig:cmp_sprint1}
\end{figure}

La struttura delle classi del modello di dominio introdotta nello Sprint 1 è delineata nella Figura~\ref{fig:class_sprint1}.

\begin{figure}[H]
    \centering
    \begin{tikzpicture}[
        node distance=0.8cm and 2.6cm,
        classbox/.style={
            draw=black!80, fill=white, thick,
            rectangle split, rectangle split parts=3,
            rounded corners=2pt, align=left, font=\scriptsize, inner sep=4pt
        }
    ]
        % Building
        \node[classbox] (building) {
            \textbf{Building}
            \nodepart{second}
            - id: String\\
            - name: String\\
            - address: String\\
            - latitude: Double\\
            - longitude: Double\\
            - surfaceM2: Double\\
            - annualEnergyConsumptionKwh: Double\\
            - annualEnergyCostEur: Double\\
            - annualCo2Kg: Double\\
            - provenance: DataProvenance
            \nodepart{third}
            + getId(): String\\
            + getName(): String
        };

        % InterventionOption
        \node[classbox, right=2.6cm of building] (option) {
            \textbf{InterventionOption}
            \nodepart{second}
            - id: String\\
            - buildingId: String\\
            - interventionTypeId: String\\
            - name: String\\
            - investmentCostCents: long\\
            - annualEnergySavingScaled: long\\
            - annualEconomicSavingCents: long\\
            - annualCo2ReductionScaled: long\\
            - provenance: DataProvenance
            \nodepart{third}
            + getId(): String\\
            + getInvestmentCostCents(): long
        };

        % InterventionType
        \node[classbox, below=1.2cm of option] (type) {
            \textbf{InterventionType}
            \nodepart{second}
            - id: String\\
            - name: String\\
            - category: String\\
            - description: String
            \nodepart{third}
            + getId(): String\\
            + getName(): String
        };

        % Connectors with clear non-overlapping labels
        \draw[-Latex, thick] (option.west) -- 
            node[above=3pt, font=\scriptsize] {1..* appartiene a} 
            node[below=3pt, font=\scriptsize] {1} (building.east);
        \draw[-Latex, thick] (option.south) -- 
            node[right=4pt, font=\scriptsize] {1..* tipologia 1} (type.north);
    \end{tikzpicture}
    \caption{Class Diagram del dominio introdotto nello Sprint 1.}
    \label{fig:class_sprint1}
\end{figure}

\subsubsection{Test Effettuati ed Evidenze di Collaudo}
\begin{itemize}[leftmargin=1.4em, itemsep=0.2em]
    \item \textbf{Test di Coerenza del Dataset (\texttt{DemoDatasetCoherenceTest}):} Verifica automatica sul dataset caricato da \texttt{DemoDataLoader}: assenza di valori nulli su campi obbligatori, correttezza dei vincoli di unicità degli identificativi e corrispondenza referenziale degli edifici (100 edifici, 300 interventi, 6 tipologie caricate con successo).
    \item \textbf{Test di Integrazione Controller REST:} Verifica delle risposte HTTP 200 su \texttt{/api/buildings}, \texttt{/api/interventions} e \texttt{/api/map/buildings}.
    \item \textbf{Test Frontend Vitest:} Collaudo del rendering della shell e delle viste catalogo (\texttt{BuildingsPage.test.tsx}, \texttt{InterventionsPage.test.tsx}, \texttt{MapPage.test.tsx}).
\end{itemize}

\subsubsection{Retrospettiva dello Sprint 1}
\begin{itemize}[leftmargin=1.4em, itemsep=0.2em]
    \item \textbf{Cosa ha funzionato:} La configurazione di Flyway e Spring Data JPA ha consentito un bootstrap immediato e riproducibile della persistenza. L'integrazione di Leaflet su coordinate reali del Comune di Trento ha reso immediatamente tangibile il valore della visualizzazione cartografica.
    \item \textbf{Criticità riscontrate:} I tile cartografici standard di fornitori terzi applicavano watermark invasivi e presentavano problemi di resa in modalità scura; la criticità è stata superata adottando i tile OpenStreetMap filtrati mediante token CSS nativi.
    \item \textbf{Azioni di miglioramento per lo sprint successivo:} Formalizzare la gestione dei valori economici in centesimi di euro (\texttt{long}) per prevenire imprecisioni di arrotondamento a virgola mobile durante l'ottimizzazione matematica.
\end{itemize}

% -------------------------------------------------------------
% SPRINT 2
% -------------------------------------------------------------
\subsection{Sprint 2: Core di Ottimizzazione Esatta CP-SAT, Vincoli e Spiegazioni}

\subsubsection{Obiettivo dello Sprint e User Story Affrontate}
Lo Sprint 2 ha realizzato il nucleo matematico ed euristico di TEPO, introducendo il solutore esatto di programmazione a vincoli, la gestione dei vincoli logici booleani e il generatore di spiegazioni deterministiche:
\begin{itemize}[leftmargin=1.5em, itemsep=0.2em]
    \item \textbf{US-03 (Configurazione Guidata Scenario):} Creazione guidata dello scenario con assegnazione del budget, scelta dell'obiettivo prioritario e applicazione di vincoli tecnici (incompatibilità tra interventi, prerequisiti tecnologici e numero massimo di interventi per tipologia).
    \item \textbf{US-04 (Risoluzione Portafoglio Ottimo con OR-Tools CP-SAT):} Calcolo deterministico della combinazione ottima di interventi mediante programmazione matematica esatta con tie-break canonico a 4 passi.
    \item \textbf{US-05 (Motore di Spiegazione Deterministico):} Produzione di motivazioni formalizzate per ogni intervento (incluso o escluso), evidenziando barriere locali di budget o incompatibilità con gli interventi selezionati.
    \item \textbf{US-14 (Audit Trail e Hash Canonico del Problema):} Calcolo dell'hash SHA-256 canonico dello scenario e registrazione persistente dell'esecuzione dell'ottimizzazione per futura ricostruibilità amministrativa.
\end{itemize}

\subsubsection{Piano di Implementazione Tecnico}
\begin{itemize}[leftmargin=1.4em, itemsep=0.2em]
    \item \textbf{Database:} Creazione della migrazione Flyway \texttt{V2\_\_optimization\_runs.sql}, introducendo la tabella \texttt{optimization\_runs} per memorizzare le soluzioni ottime, il problem hash crittografico, i reason code e i tempi di computazione.
    \item \textbf{Backend:} Integrazione dei binari nativi Google OR-Tools 9.11 tramite JNI. Implementazione del solutore esatto \texttt{OrToolsSolver} che traduce il problema in variabili decisionali binarie $x_i \in \{0, 1\}$. Implementazione della procedura di tie-break a 4 passi:
    \begin{enumerate}[label=\alph*), itemsep=0.1em]
        \item Massimizzazione dell'obiettivo primario $Z_1 = \sum_{i} w_i x_i$ con vincolo $\sum_i c_i x_i \le B$;
        \item Con $Z_1$ vincolato al valore ottimo $Z_1^*$, minimizzazione dell'investimento totale $Z_2 = \sum_i c_i x_i$;
        \item Con $Z_1^*$ e $Z_2^*$ vincolati, minimizzazione del numero totale di interventi realizzati $Z_3 = \sum_i x_i$;
        \item Spareggio lessicografico rigoroso sugli identificativi univoci delle opzioni candidate.
    \end{enumerate}
    Sviluppo del modulo \texttt{ExplanationService} per la verifica delle barriere locali di saturazione e del \texttt{ProblemHasher} per la generazione di hash canonici deterministici.
    \item \textbf{Frontend:} Realizzazione della vista di composizione scenario \texttt{ScenarioBuilder.tsx} e della schermata di dettaglio esiti \texttt{ResultPage.tsx}, completa di cruscotto KPI, prospetto di allocazione del budget ed esploratore motivazioni.
\end{itemize}

\subsubsection{Evoluzione Architetturale dello Sprint 2}
La Figura~\ref{fig:cmp_sprint2} illustra l'integrazione del motore di ottimizzazione nativo e del servizio di spiegazione all'interno dell'Application Layer.

\begin{figure}[H]
    \centering
    \begin{tikzpicture}[
        node distance=0.8cm and 0.8cm,
        comp/.style={
            draw=black!80, fill=tepoblue!10, rounded corners=3pt, thick,
            align=center, font=\scriptsize\bfseries, minimum height=0.9cm, text width=3.3cm
        },
        solvercomp/.style={
            draw=orange!80, fill=orange!15, rounded corners=3pt, thick,
            align=center, font=\scriptsize\bfseries, minimum height=0.9cm, text width=3.4cm
        },
        db/.style={
            draw=black!80, fill=yellow!15, rounded corners=4pt, thick,
            align=center, font=\scriptsize\bfseries, minimum height=1.0cm, text width=5.0cm
        },
        box/.style={
            draw=tepoblue!60, fill=tepoblue!3, rounded corners=4pt, dashed, inner sep=7pt
        },
        arrow/.style={-Latex, thick, draw=black!80}
    ]
        % Frontend Tier
        \node[comp] (views2) at (-2.4, 0) {ScenarioBuilder /\\ResultPage};
        \node[comp] (api2) at (2.4, 0) {ApiClient\\(TypeScript)};
        \begin{scope}[on background layer]
            \node[box, fit=(views2) (api2), label=above:{\scriptsize\textbf{Frontend SPA}}] (b_fe2) {};
        \end{scope}

        % Backend Controllers & Orchestration
        \node[comp] (ctrl2) at (-2.6, -1.8) {ScenarioController\\(REST Endpoint)};
        \node[comp] (serv2) at (2.6, -1.8) {OptimizationService\\(Orchestrator)};

        % Solver and Analysis Sub-components in a balanced row (zero overlap)
        \node[comp, text width=3.3cm] (hasher) at (-4.2, -3.7) {ProblemHasher\\(SHA-256)};
        \node[solvercomp, text width=3.3cm] (solver) at (0, -3.7) {OrToolsSolver\\(C++ Engine JNI)};
        \node[comp, text width=3.3cm] (expl) at (4.2, -3.7) {ExplanationService\\(Local Barriers)};

        \begin{scope}[on background layer]
            \node[box, fit=(ctrl2) (serv2) (hasher) (solver) (expl), label=above:{\scriptsize\textbf{Backend Application \& Optimization Tier}}] (b_be2) {};
        \end{scope}

        % DB Tier centered below
        \node[db] (db2) at (0, -5.6) {PostgreSQL 16\\(Flyway V1 + V2)};

        % Connectors - Zero crossing
        \draw[arrow] (views2) -- (api2);
        \draw[arrow] (api2.south) |- (ctrl2.east);
        \draw[arrow] (ctrl2) -- (serv2);

        \draw[arrow] (serv2.south) -| (hasher.north);
        \draw[arrow] (serv2.south) -- (solver.north);
        \draw[arrow] (serv2.south) -| (expl.north);

        \draw[arrow] (solver.south) -- node[right, font=\tiny] {Audit Runs} (db2.north);
        \draw[arrow] (ctrl2.west) -- node[above, font=\tiny] {Lettura Catalogo} ++(-2.4cm, 0) |- (db2.west);
    \end{tikzpicture}
    \caption{Component Diagram allo stato dello Sprint 2.}
    \label{fig:cmp_sprint2}
\end{figure}

Il raffinamento del Class Diagram con l'introduzione delle strutture matematiche e dei contratti di risoluzione è illustrato in Figura~\ref{fig:class_sprint2}.

\begin{figure}[H]
    \centering
    \begin{tikzpicture}[
        node distance=0.8cm and 1cm,
        classbox/.style={draw=black!80, fill=white, thick, rectangle split, rectangle split parts=3, rounded corners=2pt, align=left, font=\scriptsize, inner sep=4pt}
    ]
        % OptimizationProblem
        \node[classbox] (prob) {
            \textbf{OptimizationProblem}
            \nodepart{second}
            - scenarioId: String\\
            - budgetCents: long\\
            - objective: Objective\\
            - candidateOptions: List<CandidateOption>\\
            - constraints: List<ScenarioConstraint>
            \nodepart{third}
            + candidateOptions(): List<CandidateOption>\\
            + budgetCents(): long
        };

        % OptimizationResult
        \node[classbox, right=1.2cm of prob] (res) {
            \textbf{OptimizationResult}
            \nodepart{second}
            - runId: String\\
            - status: SolverStatus\\
            - totalInvestmentCents: long\\
            - objectiveValue: long\\
            - selectedOptionIds: List<String>\\
            - explanations: List<Explanation>\\
            - problemHash: String
            \nodepart{third}
            + isOptimal(): boolean\\
            + selectedOptionIds(): List<String>
        };

        % Solver Interface & Impl
        \node[classbox, below=1cm of prob] (solver_int) {
            \textbf{<<interface>> OptimizationSolver}
            \nodepart{second}
            \nodepart{third}
            + solve(p: OptimizationProblem): OptimizationResult
        };

        \node[classbox, below=0.6cm of solver_int] (ortools) {
            \textbf{OrToolsSolver}
            \nodepart{second}
            - solverConfiguration: SolverConfiguration
            \nodepart{third}
            + solvePrimary(...): StepResult\\
            + solveMinCost(...): StepResult\\
            + solveMinCount(...): StepResult\\
            + solveLexicographic(...): List<String>
        };

        % Connectors
        \draw[-Latex, thick] (solver_int) -- (prob);
        \draw[-Latex, thick] (solver_int) -- (res);
        \draw[-Latex, dashed, thick] (ortools) -- (solver_int);
    \end{tikzpicture}
    \caption{Class Diagram del modulo di ottimizzazione (Sprint 2).}
    \label{fig:class_sprint2}
\end{figure}

\subsubsection{Test Effettuati ed Evidenze di Collaudo}
Lo Sprint 2 è stato convalidato da una vasta suite di test algoritmici mirati:
\begin{itemize}[leftmargin=1.4em, itemsep=0.2em]
    \item \textbf{Test di Parità Solutore (\texttt{SolverParityTest}, \texttt{SolverRandomizedParityTest}):} Verifica comparativa tra \texttt{OrToolsSolver} e un risolutore esaustivo brute-force (\texttt{BruteForceSolver}) su 100 istanze generate casualmente. L'uguaglianza rigorosa della funzione obiettivo e delle selezioni ha raggiunto il 100\% di conformità.
    \item \textbf{Test di Confronto Baseline Greedy (\texttt{GreedyBaselineComparisonTest}):} Dimostrazione analitica della superiorità del solver esatto rispetto all'algoritmo ingoloso euristico, con scarti fino al 18\% su istanze con vincoli di incompatibilità.
    \item \textbf{Test di Regressione Tie-Break (\texttt{OrToolsSolverRegressionTest}):} Verifica dell'indipendenza dell'output rispetto alla permutazione dell'ordine dei dati di input.
    \item \textbf{Test di Spiegazione (\texttt{ExplanationReasonCodeTest}):} Validazione dell'assegnazione corretta dei codici di esclusione (\texttt{EXCLUDED\_BUDGET\_BARRIER}, \texttt{EXCLUDED\_INCOMPATIBILITY}, ecc.).
\end{itemize}

\subsubsection{Retrospettiva dello Sprint 2}
\begin{itemize}[leftmargin=1.4em, itemsep=0.2em]
    \item \textbf{Cosa ha funzionato:} Le prestazioni di Google OR-Tools si sono rivelate straordinarie: problemi a 300 variabili binarie con vincoli logici multipli vengono risolti all'ottimo globale in meno di 150 millisecondi.
    \item \textbf{Criticità riscontrate:} La procedura di spareggio lessicografico necessitava della ricostruzione atomica del modello CP-SAT ad ogni fissaggio di variabile binaria per evitare lo stallo dei vincoli interni di OR-Tools.
    \item \textbf{Azioni di miglioramento per lo sprint successivo:} Fornire un supporto strutturato al confronto visuale affiancato di scenari per consentire ai decisori di valutare l'impatto di variazioni di budget.
\end{itemize}

% -------------------------------------------------------------
% SPRINT 3
% -------------------------------------------------------------
\subsection{Sprint 3: Analisi What-If, Sensibilità sul Budget e Importazione CSV}

\subsubsection{Obiettivo dello Sprint e User Story Affrontate}
Lo Sprint 3 ha introdotto gli strumenti decisionali di livello avanzato per l'esplorazione parametrica della stabilità delle soluzioni e l'ingestione controllata di nuovi dati:
\begin{itemize}[leftmargin=1.5em, itemsep=0.2em]
    \item \textbf{US-06 (Confronto Scenari What-If A/B):} Memorizzazione di due configurazioni decisionali distinte in slot di sessione denominati ``Slot A'' e ``Slot B'', con analisi differenziale affiancata e calcolo automatico dei delta ($\Delta = \text{B} - \text{A}$) su spesa, risparmio kWh, risparmio economico ed emissioni di CO\textsubscript{2}.
    \item \textbf{US-07 (Analisi di Sensibilità sul Budget):} Esecuzione di uno sweep parametrico deterministico variando il budget lungo una griglia uniforme configurabile da 3 a 51 punti, classificando gli interventi in \emph{stabili} (sempre inclusi o sempre esclusi) oppure \emph{instabili/oscillanti}.
    \item \textbf{US-10 (Importazione Guidata Dati da CSV):} Procedura di caricamento in tre fasi (upload e parsing tollerante, validazione a secco atomica con rilevazione errori e commit transazionale).
\end{itemize}

\subsubsection{Piano di Implementazione Tecnico}
\begin{itemize}[leftmargin=1.4em, itemsep=0.2em]
    \item \textbf{Backend:} Implementazione del servizio \texttt{SensitivityService} per l'orchestrazione dello sweep lineare sul budget con invocazione sequenziale del solver e generazione delle matrici di frequenza di inclusione. Sviluppo del modulo \texttt{ImportController} dotato di parser CSV tollerante (gestione automatica di separatori virgola/punto e virgola, apici e codifiche UTF-8/ISO-8859) e meccanismo di rollback completo in caso di violazione dell'integrità referenziale.
    \item \textbf{Frontend:} Sviluppo dello store client di sessione \texttt{scenarioStore.ts} in LocalStorage; creazione delle viste operative \texttt{ComparePage.tsx} (tabelle comparative con delta cromatici), \texttt{SensitivityPage.tsx} (grafici SVG nativi per curve di efficienza e heatmap di stabilità) e \texttt{ImportPage.tsx} (wizard con anteprima errori riga per riga e statistiche di caricamento).
\end{itemize}

\subsubsection{Evoluzione Architetturale dello Sprint 3}
La Figura~\ref{fig:cmp_sprint3} illustra l'inclusione dei sottosistemi di sensibilità parametrica e di ingestione dei dati esterni.

\begin{figure}[H]
    \centering
    \begin{tikzpicture}[
        node distance=0.8cm and 1cm,
        comp/.style={draw=black!80, fill=tepoblue!10, rounded corners=3pt, thick, align=center, font=\scriptsize\bfseries, minimum height=0.9cm, minimum width=3.2cm},
        servcomp/.style={draw=tepogreen!80, fill=tepogreen!15, rounded corners=3pt, thick, align=center, font=\scriptsize\bfseries, minimum height=0.9cm, minimum width=3.4cm},
        box/.style={draw=tepoblue!60, fill=tepoblue!3, rounded corners=4pt, dashed, inner sep=7pt},
        arrow/.style={-Latex, thick, draw=black!80}
    ]
        % Frontend
        \node[comp] (views3) {ComparePage / SensitivityPage /\\ImportPage};
        \node[comp, right=0.8cm of views3] (store) {scenarioStore (Client A/B)};
        \begin{scope}[on background layer]
            \node[box, fit=(views3) (store), label=above:{\scriptsize\textbf{Frontend State \& Views (Sprint 3)}}] (b_fe3) {};
        \end{scope}

        % Backend
        \node[comp, below=1cm of views3] (ctrl3) {SensitivityController /\\ImportController};
        \node[servcomp, below=0.8cm of ctrl3] (sens_serv) {SensitivityService\\(Budget Sweep Engine)};
        \node[servcomp, right=0.8cm of sens_serv] (opt_serv) {OptimizationService};
        \begin{scope}[on background layer]
            \node[box, fit=(ctrl3) (sens_serv) (opt_serv), label=above:{\scriptsize\textbf{Backend Application Layer}}] (b_be3) {};
        \end{scope}

        % Connectors
        \draw[arrow] (views3) -- (store);
        \draw[arrow] (store) |- (ctrl3);
        \draw[arrow] (ctrl3) -- (sens_serv);
        \draw[arrow] (sens_serv) -- (opt_serv);
    \end{tikzpicture}
    \caption{Component Diagram allo stato dello Sprint 3.}
    \label{fig:cmp_sprint3}
\end{figure}

\subsubsection{Test Effettuati ed Evidenze di Collaudo}
\begin{itemize}[leftmargin=1.4em, itemsep=0.2em]
    \item \textbf{Test Sweep di Sensibilità (\texttt{SensitivityServiceTest}):} Verifica della correttezza della griglia di campionamento (punti equidistanti tra budget minimo e massimo), monotonicità non decrescente del valore obiettivo e corretta ripartizione degli interventi in categorie di stabilità.
    \item \textbf{Test Importazione CSV (\texttt{ImportControllerTest}):} Collaudo della validazione atomica: rigetto di file con colonne mancanti, righe con identificativi duplicati o chiavi esterne inesistenti senza alcuna persistenza sul database; accettazione e caricamento con rollback opzionale tramite l'endpoint \texttt{/api/import/reset-demo}.
    \item \textbf{Test Frontend Vitest:} Verifica dell'interazione nelle pagine \texttt{ComparePage.test.tsx}, \texttt{SensitivityPage.test.tsx} e \texttt{ImportPage.test.tsx}.
\end{itemize}

\subsubsection{Retrospettiva dello Sprint 3}
\begin{itemize}[leftmargin=1.4em, itemsep=0.2em]
    \item \textbf{Cosa ha funzionato:} Lo store di sessione client ha disaccoppiato la gestione degli slot A/B dal backend, consentendo all'utente di alternare, scambiare e confrontare configurazioni decisionali con latenza zero.
    \item \textbf{Criticità riscontrate:} L'esecuzione dello sweep su 51 punti di budget generava un carico temporaneo significativo se eseguita senza vincoli di timeout; è stato introdotto un limite massimo configurabile sul tempo di risoluzione per singolo punto.
    \item \textbf{Azioni di miglioramento per lo sprint successivo:} Completare l'assistente AI con architettura di fallback resiliente e perfezionare l'accessibilità visiva e l'infrastruttura di rilascio in cloud.
\end{itemize}

% -------------------------------------------------------------
% SPRINT 4
% -------------------------------------------------------------
\subsection{Sprint 4: Assistente AI Sovrintendente, Polishing e Rilascio Cloud}

\subsubsection{Obiettivo dello Sprint e User Story Affrontate}
L'ultimo sprint ha integrato il modulo di intelligenza artificiale opzionale secondo i canoni della piena responsabilità umana, perfezionando l'accessibilità del sistema e finalizzando la pipeline di deployment automatizzato:
\begin{itemize}[leftmargin=1.5em, itemsep=0.2em]
    \item \textbf{US-11 (Assistente AI Human-in-the-Loop):} Modulo di interpretazione delle richieste in linguaggio naturale dell'operatore tecnico (es. ``Massimizza il risparmio per le scuole con un budget di 500.000 euro'') per proporre parametri di scenario strutturati senza alcuna scrittura diretta sui dati.
    \item \textbf{US-12 (Fallback Rule-Based Locale Deterministico):} Meccanismo di failover trasparente che, in assenza di connessione o in mancanza di chiavi API per il provider LLM esterno, elabora la richiesta mediante un motore a regole euristico locale, segnalando lo stato all'utente.
    \item \textbf{US-13 (Accessibilità WCAG 2.1 AA e Doppio Registro Visivo):} Pieno supporto per ipovedenti, contrasto elevato, navigabilità integrale da tastiera e supporto a due temi operativi (\emph{Relazione} per lettura/proiezione diurna e \emph{Sala Operativa} per sessioni analitiche prolungate).
    \item \textbf{US-15 (Esportazione e Stampa della Relazione Tecnica):} Foglio di stile CSS per la stampa che consente di generare un report documentale pulito e stampabile degli esiti decisionali.
\end{itemize}

\subsubsection{Piano di Implementazione Tecnico}
\begin{itemize}[leftmargin=1.4em, itemsep=0.2em]
    \item \textbf{Backend:} Implementazione dell'interfaccia \texttt{AiAssistant} e delle due classi concrete \texttt{LlmAiAssistant} (integrazione modello linguistico esterno con formattazione JSON controllata) e \texttt{RuleBasedAiAssistant} (motore deterministico basato su espressioni regolari e corrispondenza tassonomica). Implementazione di \texttt{AiController} per esporre lo stato operativo (\texttt{/api/ai/status}) e la generazione delle proposte (\texttt{/api/ai/propose-scenario}). Sviluppo del componente \texttt{DatabaseUrlEnvironmentPostProcessor} per la risoluzione trasparente delle variabili di connessione del cloud provider Render.
    \item \textbf{Frontend:} Realizzazione del pannello estensibile \texttt{AiAssistantPanel.tsx} nello Scenario Builder; introduzione della tavolozza comandi rapida \texttt{CommandPalette.tsx} (\texttt{Ctrl+K}); ottimizzazione del bundle statico per Nginx e stili di stampa (@media print).
    \item \textbf{Deployment \& Cloud:} Configurazione del blueprint Infrastructure-as-Code \texttt{render.yaml}, validazione dei container multi-stage Docker e collaudo dell'ambiente pubblico su \url{https://tepo.onrender.com}.
\end{itemize}

\subsubsection{Evoluzione Architetturale dello Sprint 4}
La Figura~\ref{fig:cmp_sprint4} illustra l'architettura integrata completata dello Sprint 4, evidenziando il sottosistema AI ibrido e l'integrazione di deployment cloud.

\begin{figure}[H]
    \centering
    \begin{tikzpicture}[
        node distance=0.8cm and 1.2cm,
        comp/.style={
            draw=black!80, fill=tepoblue!10, rounded corners=3pt, thick,
            align=center, font=\scriptsize\bfseries, minimum height=0.9cm, text width=4.5cm
        },
        aicomp/.style={
            draw=purple!80, fill=purple!15, rounded corners=3pt, thick,
            align=center, font=\scriptsize\bfseries, minimum height=0.9cm, text width=4.5cm
        },
        cloudcomp/.style={
            draw=tepogreen!80, fill=tepogreen!15, rounded corners=3pt, thick,
            align=center, font=\scriptsize\bfseries, minimum height=0.9cm, text width=4.5cm
        },
        box/.style={
            draw=tepoblue!60, fill=tepoblue!3, rounded corners=4pt, dashed, inner sep=6pt
        },
        arrow/.style={-Latex, thick, draw=black!80}
    ]
        % Frontend & Reverse Proxy
        \node[comp, text width=4.8cm] (ui_s4) at (-3.4, 0) {React 18 SPA\\(AiAssistantPanel + Shell)};
        \node[comp, text width=4.8cm] (nginx_s4) at (3.4, 0) {Nginx Reverse Proxy\\(Docker Web Service)};

        % Backend
        \node[comp, text width=8.8cm] (api_s4) at (0, -1.8) {
            Spring Boot 3.3 Application Tier\\
            \tiny REST Controllers $\cdot$ OptimizationService $\cdot$ DatabaseUrlPostProcessor
        };

        % Subsystems side-by-side below (Zero overlap)
        \node[aicomp, text width=5.6cm] (ai_engine) at (-3.5, -3.8) {
            AiAssistant Subsystem\\
            \tiny LLM Client + Rule-Based Fallback Locale
        };
        \node[cloudcomp, text width=5.6cm] (cloud_s4) at (3.5, -3.8) {
            Render PaaS Cloud Environment\\
            \tiny Managed PostgreSQL 16 + Auto-Deploy Webhook
        };

        % Connectors - Clean orthogonal routing
        \draw[arrow] (ui_s4) -- node[above, font=\tiny] {HTTPS} (nginx_s4);
        \draw[arrow] (nginx_s4.south) -- node[right, font=\tiny] {Proxy pass (:8080)} ++(0, -0.6cm) -| (api_s4.north);
        \draw[arrow] (api_s4.south) ++(-2.4cm, 0) -- ++(0, -0.5cm) -| (ai_engine.north);
        \draw[arrow] (api_s4.south) ++(2.4cm, 0) -- ++(0, -0.5cm) -| (cloud_s4.north);
    \end{tikzpicture}
    \caption{Component Diagram integrato allo stato dello Sprint 4.}
    \label{fig:cmp_sprint4}
\end{figure}

\subsubsection{Test Effettuati ed Evidenze di Collaudo}
\begin{itemize}[leftmargin=1.4em, itemsep=0.2em]
    \item \textbf{Test Assistente AI (\texttt{LlmAiAssistantTest}, \texttt{RuleBasedAiAssistantTest}):} Verifica della resilienza del fallback rule-based a fronte di timeout del provider esterno o chiavi non configurate; conformità dello schema JSON restituito con i vincoli di dominio.
    \item \textbf{Test di Configurazione Cloud (\texttt{DatabaseUrlEnvironmentPostProcessorTest}):} Verifica della corretta traduzione dell'URI standard `postgres://` fornito da Render nella stringa JDBC conforme `jdbc:postgresql://` con disattivazione forzata di credenziali concorrenti.
    \item \textbf{Collaudo di Accessibilità WCAG 2.1 AA:} Verifica dei rapporti di contrasto cromatico ($> 4.5:1$ per testo normale, $> 3:1$ per componenti grafici) sia nel tema Relazione che nel tema Sala Operativa.
    \item \textbf{Collaudo Live in Cloud:} Esecuzione end-to-end dello scenario di test sull'istanza di produzione ospitata su Render con superamento di tutti i controlli di salute (\texttt{healthcheck} HTTP 200 su \texttt{/api/scenarios/ping}).
\end{itemize}

\subsubsection{Retrospettiva dello Sprint 4}
\begin{itemize}[leftmargin=1.4em, itemsep=0.2em]
    \item \textbf{Cosa ha funzionato:} La separazione netta tra LLM esterno e fallback a regole ha eliminato ogni punto singolo di fallimento, consentendo al sistema di operare fluidamente anche in totale assenza di connettività internet.
    \item \textbf{Criticità riscontrate:} La modalità free-tier di Render pone il servizio in stato di sospensione dopo 15 minuti di inattività, provocando un avvio a freddo (cold start) di circa 30-45 secondi alla prima richiesta; la situazione è stata documentata nel manuale utente.
    \item \textbf{Sintesi del rilascio finale:} Il team ha completato con successo il 100\% delle User Story pianificate (87 Story Points totali), raggiungendo una copertura di test automatici completa e rilasciando un prodotto stabile, verificabile e pienamente conforme alle specifiche.
\end{itemize}
